\section{Scan controls}
\label{sec:scan}

\subsection{Scan parameters}

\begin{center}
\begin{tabularx}{\textwidth}{@{}lL@{}}
\toprule
Field & Meaning \\
\midrule
\app{Scan type}          & Sawtooth (repeats up), Triangle (up and back) or Single scan \\
\app{Bunches per channel} & bunches taken at each setpoint; \app{Seconds per channel} in DC mode \\
\app{Stepsize (V)}       & the step between setpoints \\
\app{From (V)}, \app{To (V)} & up to seven ranges, each ticked to be scanned; a point
                           shared by two ranges is visited once \\
\app{Experiment name}    & recorded with every run and in the ELOG \\
\app{Interaction region} & CEC or LCR: the cooler offset differs between them \\
\app{PMT HV}             & the PMTs' high voltage, recorded where EPICS does not give it \\
\app{Reference run}      & marks a run taken to measure the beam's offset from the cooler \\
\bottomrule
\end{tabularx}
\end{center}

Every voltage is held to the platform's limit (by default the Trek 609E-6's $\pm$4\,kV). A
range that is not a whole number of steps ends on its last whole step, never past its stop.

\subsection{Buncher controls}

\app{Bunched mode} takes one bunch per \app{Dwell time}: the buncher is released for
\app{Bunch release length}, and the tagger's window opens \app{DAQ TOF trigger delay} after
the release and stays open from \app{DAQ gate starts} to \app{DAQ gate stops}, inside the
tagger's 8.39\,ms range. \app{DC mode} opens a window every millisecond; its gate stays
settable, since DC hits come through the tagger too. A timing that does not fit its cycle
(a release or delay past its end, a window still open at the next trigger) is refused with
a sentence before anything starts.

\subsection{Starting and stopping}

\app{Start scan} calibrates the scan's own setpoints (Section~\ref{sec:calibration}); only
if the gate passes are the run file written and the scan started. A refused calibration
writes nothing, and its run number is given to the next attempt. \app{Stop after scan}
finishes the scan in progress; \app{Stop now!} ends within a bunch.

At each step the tagger's window is shut, the DAC written, and the window held shut for as
long as the settle model says the jump takes; every step then takes its full count of
bunches at a settled voltage, while the buncher and cooler run on untouched. A tagger that
sees no trigger for three cycles (two seconds at least) is reported on the status bar, and
the scan holds its setpoint until triggers return.

\subsection{Scan queue}

The second page of the parameters box lays scans out ahead, one row and one run file each:
ranges (\code{20:55, 60:80}), step, bunches (seconds in DC), order (Sawtooth or Triangle)
and number of scans; timing, experiment and view are the parameters page's. Each row is
checked as it is typed (a bad cell turns red) and the readout gives each row's setpoints and
time, or its fault, and when the queue ends. \app{Add current} adds the parameters as they
stand, \app{Show} puts a row back on them, rows paste from a spreadsheet, Delete clears
cells.

\app{Automatic} starts each scan as the last ends, so a queue can run overnight;
\app{Manual} asks before each scan after the first. Anything but a clean end (a calibration
refused, a file not written, a scan stopped by hand) stops the queue with the platform at
rest, its row saying why. \app{Stop queue after this scan} lets the running scan finish.
Rows and states survive a crash; a restarted queue takes up at the first scan not done.

\subsection{Plots}

The hyperfine-structure spectrum, the time-of-flight (TOF) spectrum and, with \app{Show
time-of-flight map}, the TOF against scan voltage. Ctrl-drag (Cmd-drag on macOS) on the TOF
spectrum sets the gate, a double-click removes it; the spectrum and map count only hits
inside it. \app{Bin size} and the \app{Channel} boxes choose the TOF binning and the PMTs.
The view is held from Start to the scan's last event and recorded with the run.

Every plot handles alike: drag to pan, right-drag to stretch, the wheel to zoom (over an
axis, that axis alone), Shift-drag to zoom to a box. A plot handled by hand stops following
the data until its \raisebox{-0.1ex}{$\circlearrowleft$} button, a double-click or
\menu{View > Autoscale all plots} (Ctrl+0) returns it. Resting the cursor on a title says
what the plot shows and how it is calculated.

\subsection{ELOG, quick tune and the experiment information}

\app{Send additional text} appends to the open ELOG entry, \app{Start new entry} begins the
next, \app{Send HFS screen} attaches the spectrum. Each scan posts its parameters and
calibration; a simulated rack posts nothing.

\app{Enable triggers} runs the triggers without a scan; \app{Enable voltage} puts
\app{Set voltage} on the platform, set by value or by the slider across its range, starting
from the voltage already there. Spin boxes and the slider take the mouse wheel only once
clicked, so scrolling past them never moves the platform. Below, the cooler and readback
voltages, the wavenumber, run and scan numbers, the buncher mode, and \app{Shown}: the run
every tab is showing, live, saved or opened.
